% Préambule APMEP - Annales
% Auteurs
% tapuscrit : Denis Vergès
% relecture : François Kriegk
% chiffres TikZ : François Kriegk
% chiffres pstricks : Denis Vergès
% modifications :
% Remerciements :
% version 3 : 30/06/2026
%
\documentclass[12pt,a4paper,french]{article}

% des paquets
%%% codage et police utilisée
\usepackage[T1]{fontenc}
\RequirePackage{iftex}
\iftutex % moteur utf-8 lualatex, xelatex
	\usepackage{fourier-otf} % remplace fourier pour moteur utf8 (LuaLaTeX, XeLaTeX)
	% Alternatives nécéssites (LuaLaTeX, XeLaTeX)
    \usepackage{hyperref}
	%\usepackage{fontspec}% pour utiliser les polices système
\else % pour les autres moteurs latex, pdflatex
	\usepackage[utf8]{inputenc}
	\usepackage{fourier}
	\usepackage{amsmath,amssymb}
    \usepackage[dvips]{hyperref}
    \renewcommand{\ttdefault}{lmtt}
\fi
\usepackage[normalem]{ulem}
\usepackage{pifont}
\usepackage{lscape}
% tableaux
\usepackage{multicol}
\usepackage{multirow}
\usepackage{tabularx}
\usepackage{tabularray}
% texte et composition
\usepackage{textcomp}
\usepackage{babel}
\usepackage[np]{numprint}
%
\usepackage{fancyhdr}
\usepackage{makeidx}
\makeindex
% pour le symbole €
\usepackage{eurosym}
\DeclareRobustCommand{\officialeuro}{%
  \ifmmode\expandafter\text\fi
  {\fontencoding{U}\fontfamily{eurosym}\selectfont e}}
\AtBeginDocument{\renewcommand{\euro}{\officialeuro}}
% couleurs étendues
\usepackage{xcolor}
% pour les listes
\usepackage{enumitem}
% Enumérations réglées par défaut
\setlist[itemize]{label=\textbullet,leftmargin=5mm,labelwidth=3mm,labelsep=2mm,align=left}
\setlist[enumerate]{labelindent=0pt, leftmargin=5mm, labelsep=2mm, labelwidth=3mm,labelsep=2mm, align= left}
\setlist[enumerate,1]{label = \textbf{\arabic*.}}
\setlist[enumerate,2]{label = \textbf{\alph*.}}
\setlist[enumerate,3]{label = \textbf{\roman*.}}
% pour l'ajustement des boîtes (chiffres…)
\usepackage{adjustbox}
% pour ne pas commencer un exercice en fin de page
\usepackage{needspace}
%
% Taille de la zone d'écriture
\usepackage[left=1.5cm,right=1.5cm,top=3cm,bottom=3cm]{geometry}
\setlength{\headheight}{15pt}
\setlength{\parindent}{0mm}
%
Formatage du PDF
\hypersetup{%
    auteur du PDF = {APMEP},
    pdfsubject = {Troisième, Série générale},
    pdftitle = {Brevet des collèges -- Session 2026},
    allbordercolors = blanc,
    pdfstartview=FitH
}
% Pour scratch (code)
\usepackage{scratch3}
% Pour le code Python
\usepackage{listings}
\lstdefinestyle{python}{
  langue=Python,
  style de base=\ttfamily\small,
  style_mot_clé=\color{bleu}\bfseries,
  style de commentaire=\color{gris}\itshape,
  style de chaîne=\color{orange},
  nombres=gauche,
  style_numéro=\tiny\color{gris},
  breaklines=true,
  cadre=unique,
  couleur de fond=\color{gris!5},
  taille de l'onglet=4
}
  % d'utilisation (exemple)
  %\begin{lstlisting}[style=python]
%def seuil():
% n = 1
% p = 0,5
% tant que p < 0,6 :
% n = n + 1
% p = p*7/15 + 1/3
% retour n
%\end{lstlisting}


% Pour une visualisation de type tableau
\usepackage{pas-tableur}
%
%% pour PsTricks
\usepackage{pst-bezier}%% AVANT les autres PST
\usepackage{pstricks-add}
\usepackage{pst-func,pst-tree,pst-circ}
\usepackage{pst-eucl}
%% pour TikZ
\usepackage{tikz,pgfplots,pgf,tkz-tab,tikzpagenodes}
\usetikzlibrary{calc,patterns,decorations,arrows.meta,patterns.meta,arrows}
\pgfplotsset{compat=1.18,/pgf/number format/.cd,use comma}
\usepackage{icomma} % gestion des espaces autour des virgules
\DeclareMathSymbol{;}{\mathbin}{operators}{"3B}% pour un espacement correct du ; en mode maths
%%%%%%%%%%%%%%%%%%%%%%%%%%%%%%%%%%%%%%%%%%
% Définitions
%
% notation en mode mathématique et hors mode mathématique
\newcommand{\R}{\ensuremath{\mathbb{R}}}
\newcommand{\N}{\ensuremath{\mathbb{N}}}
\newcommand{\D}{\ensuremath{\mathbb{D}}}
\newcommand{\Z}{\ensuremath{\mathbb{Z}}}
\newcommand{\Q}{\ensuremath{\mathbb{Q}}}
\newcommand{\C}{\ensuremath{\mathbb{C}}}
%
\providecommand{\up}{\textsuperscript}
%
\newcommand{\vect}[1]{\overrightarrow{\,\mathstrut#1\,}}
\newcommand{\vectt}[1]{\overrightarrow{\,\mathstrut\text{#1}\,}}
%
\def\Oij{$\left(\text{O};\vect{\imath},\vect{\jmath}\right)$}
\def\Oijk{$\left(\text{O};\vect{\imath},\vect{\jmath},\vect{k}\right)$}
\def\Ouv{$\left(\text{O};\vect{u},\vect{v}\right)$}
%
\newcommand{\ds}{\displaystyle}
%
\renewcommand{\d}{\mathrm{\,d}}
\renewcommand{\i}{\mathrm{\,i\,}}
\newcommand{\e}{\mathrm{\,e\,}}
%%%%%%%%%%%%%%%%%%%%%%%%%%%%%%%%%%%%%%%%%%%


\begin{document}

%---------------------------------------------------------
% Réglage en tête, pied de page
%---------------------------------------------------------
\rhead{\textbf{A. P{}. ME P{}.}}
\lhead{\small Sujet}
\lfoot{\small{Asie}}
\rfoot{\small{15 juin 2026}}
\pagestyle{fancy}

\thispagestyle{vide}


%---------------------------------------------------------
% Titre
%---------------------------------------------------------
\begin{center}
{\Large \textbf{\decofourleft~ Brevet des collèges ~\decofourright\\[10pt]Sujet\\[10pt] Série générale -- Asie -- 15 juin 2026 }}
\end{center}

\marginpar{\rotatebox{90}{\textbf{A. P{}. ME P{}.}}}

\section*{Partie 1 – Automatismes – 6 points – 20 minutes}

Pour chaque question, recopier sur la copie son numéro et la réponse correspondante.

Pour cette partie, aucune justification n'est demandée.

Pour les questions à choix multiples, une seule réponse est exacte.

\subsection*{Question 1}
L'écriture scientifique du nombre \np{45310} est :

\begin{tblr}{colspec={*{4}{|X[c]}|},hlines}
	\textbf{A.}\quad $45,31 \times 10^{3}$ &
	\textbf{B.}\quad $4,531 \times 10^{4}$ &
	\textbf{C.}\quad $4,531 \times 10^{-4}$ &
	\textbf{D.}\quad $4531 \times 10^{1}$\\
\end{tblr}

\subsection*{Question 2}
Une forme développée de l'expression $(4x-3)(4x+3)$ est :

\begin{tblr}{colspec={*{4}{|X[c]}|},hlines}
	\textbf{Réponse A} & \textbf{Réponse B} & \textbf{Réponse C} & \textbf{Réponse D} \\
	4x² - 9 et 16x² + 9 et 16x² - 9 et 8x² - 6
\end{tblr}

\subsection*{Question 3}
Un pavé droit à pour dimensions : 4,5 cm de long, 4 cm de large, 10 cm de haut. Le volume de ce pavé est de :

\begin{tblr}{colspec={*{4}{|X[c]}|},hlines}
	\textbf{Réponse A} & \textbf{Réponse B} & \textbf{Réponse C} & \textbf{Réponse D} \\
	180 cm³, 170 cm³, 160,5 cm³ et 18,5 cm³
\end{tblr}

\subsection*{Question 4}
On considère les nombres suivants et on s'intéresse à leur divisibilité par 9.

\[N = \np{2025} \quad \text{et} \quad P = \np{2026}.\]

\begin{tblr}{colspec={*{4}{|X[c]}|},hlines}
	\textbf{Réponse A} & \textbf{Réponse B} & \textbf{Réponse C} & \textbf{Réponse D} \\
	\small $N$ et $P$ sont tous les deux divisibles par 9 & $N$ est divisible par 9 mais $P$ ne l'est pas & $P$ est divisible par 9 mais $N$ ne l'est pas & Aucun des deux n'est divisible par 9
\end{tblr}

\subsection*{Question 5}
Une personne a parcouru 9~km en 45~minutes.

Quelle est sa vitesse moyenne en km/h ?

\subsection*{Question 6}

\begin{minipage}{0.68\textwidth}
	Une roue de la fortune est utilisée pour faire gagner des cadeaux. La roue est divisée en 10 secteurs de tailles égales, avec les gains suivants : des stylos, des porte-clés, des casques audios ou un smartphone.

	Un joueur tourne la roue une seule fois.

	Quelle est la probabilité que le joueur gagne un casque audio ?
\end{minipage}%
\hfill
\begin{minipage}{0.28\textwidth}
	centrer
	\begin{tikzpicture}[scale=1.1]
		\foreach \i in {0,...,9} {
			\draw (0,0) -- (90+36*\i:1.6);
		}
		\dessiner (0,0) cercle (1,6);
		\node à (108:1.1) {\tiny stylo};
		\node à (72:1.1) {\tiny casque};
		\node at (36:1.1) {\tiny porte-clé};
		\node à (0:1.1) {\tiny smartphone};
		\node at (-36:1.1) {\tiny porte-clé};
		\node à (-72:1.1) {\tiny stylo};
		\node à (-108:1.1) {\tiny stylo};
		\node à (-144:1.1) {\tiny casque};
		\node at (180:1.1) {\tiny porte-clé};
		\node à (144:1.1) {\tiny stylo};
	\end{tikzpicture}
\end{minipage}

\subsection*{Question 7}
Un article coûte 60 \euro. Calculer son nouveau prix après une baisse de 10\,\%.

\subsection*{Question 8}

\begin{minipage}{0.62\textwidth}
	Quelle est la mesure de l'angle $\widehat{\text{BAC}}$ ?
\end{minipage}%
\hfill
\begin{minipage}{0.34\textwidth}
	centrer
	\begin{tikzpicture}[scale=0.8]
		\coordonnerie (C) à (0,0);
		\coordonnerie (A) à (3,0);
		\coordonnation (B) à (0,3.5);
		\dessiner (B) -- (C) -- (A) -- (B);
		\draw (0,0.3) -- (0.3,0.3) -- (0.3,0);
		\node à (-0.3,3.7) {B};
		\node[en bas à gauche] à (C) {C};
		\node[en bas à droite] à (A) {A};
		\draw (0,2.8) arc(-90:-53:0.8) node[pos=0.7, below]{$40^\circ$};
	\end{tikzpicture}
\end{minipage}

\subsection*{Question 9}
Le diagramme en barres ci-dessous donne les notes des élèves d'une classe au dernier contrôle de mathématiques.

\begin{enumerate}[label=\textbf{\alph*.}]
	\item Combien d'élèves ont participé à ce contrôle ?
	\item Quelle est la note médiane ?
\end{énumérer}

\begin{center}
	%Début Tikz
	%\begin{tikzpicture}
	%\begin{axis}[
	% ybarre,
	% largeur=15cm,
	% hauteur=6cm,
	% largeur de la barre=12pt,
	% xlabel={notes},
	% étiquette_y={effectif},
	% coordonnées x symboliques={7,8,9,10,11,12,13,14,15,16,17,18},
	% xtick=données,
	% ymin=0, ymax=6,
	% ytick={0,1,2,3,4,5,6},
	% agrandir x limites=0,05,
	% lignes d'axe=gauche,
	% nœuds proches des coordonnées,
	%]
	%\addplot coordonnées {(7,3) (8,4) (10,4) (11,5) (12,5) (15,3) (17,2) (18,1)};
	%\end{axis}
	%\end{tikzpicture}
	%Fin Tikz
	% Début PSTricks
	\psset{unit=0.9cm,arrowsize=2pt 3}
	\begin{pspicture}(-1,-1)(12.5,6.5)
		\psaxes[linewidth=1.25pt,Ox=6]{->}(0,0)(0,0)(12.5,6.5)
		\psframe(0.8,0)(1.2,3)\psframe(1.8,0)(2.2,4)\psframe(3.8,0)(4.2,4)
		\psframe(4.8,0)(5.2,5)\psframe(5.8,0)(6.2,5)\psframe(8.8,0)(9.2,3)
		\psframe(10.8,0)(11.2,2)\psframe(11.8,0)(12.2,1)
		\rput(6,-0.8){notes}\rput{90}(-0.8,3.25){effectifs}
	\end{pspicture}
	%Fin PSTricks
\end{center}

\begin{center}
	\textit{Restitution de la copie du candidat à l'issue de la parti 1}
\end{center}

\section*{Partie 2 -- Raisonnement et résolution de problèmes -- 14 points}

Dans cette partie, toutes les réponses doivent être justifiées, sauf si une indication contraire est donnée.

La clarté et la précision des raisonnements ainsi que la rédaction sont réalisées sur 2 points.

Pour chaque question, si le travail n'est pas terminé, laisser tout de même une trace de la recherche~; les essais et les démarches engagées, même non aboutis, seront pris en compte dans la notation.

\medskip

\section*{Exercice 1 \hfill 2,5 points}

Lola souhaite acheter un smartphone. Elle étudie deux propositions.

\begin{center}
	\begin{tblr}{width=0.9\linewidth,colspec={|X[7cm]| X |X[7cm]|},hline{1,Z}={1,3}{solid}}
		\textbf{Offre A} && \textbf{Offre B} \\
		Le client paie 175~\euro{} à l'achat, puis son abonnement est de 16~\euro{} par mois avec un engagement de 24 mois minimum. &&
		Le client ne paie rien à l'achat, puis l'abonnement est de 23~\euro{} par mois avec un engagement de 24 mois minimum.
	\end{tblr}
\end{center}

\begin{énumérer}
	\item Au bout de 24 mois, laquelle des deux offres est la plus intéressante ?
	\item $x$ est un nombre positif qui représente le nombre de mois. On exprime le prix de ces deux tarifs en fonction de $x$, avec les fonctions suivantes :

	\begin{itemize}
		$f(x) = 175 + 16x$
		\item $g(x) = 23x$
	\end{itemize}

	\begin{enumerate}[label={\textbf{\alph*.}}]
		\item Associer chaque fonction à l'offre correspondante (A ou B).

		\emph{Aucune justification n'est attendue.}
		\item Au bout de combien de mois payer-t-on le même prix avec ces deux offres ?
		\item Est-on encore dans la période d'engagement ?
	\end{énumérer}
\end{énumérer}

\newpage
\section*{Exercice 2 \hfill 3 points}

\emph{La figure ci-dessous n'est pas à l'échelle.}

\begin{minipage}{0.48\linewidth}
	\begin{itemize}
		\item O, A et B sont alignés.
		\item O, D et C sont alignés.
		\item OD = 8,2 cm
		\item AD $= 1,8$ cm
		\item BC $= 4,5$ cm
	\end{itemize}
\end{minipage}\hfill
\begin{minipage}{0.48\linewidth}
	% Début PSTricks
	\psset{unité=1cm}
	\begin{pspicture}(-1,-0.5)(3.6,5)
		%\psgrid
		\pspolygon(0,0)(0,5)(2.6,5)%ABCO
		\psline(0,3)(1.6,3)%AD
		\psframe(0,3)(0.3,2.7,3)
		\psframe(0,5)(0.3,4.7)
		\uput[d](0,0){O} \uput[l](0,3){A} \uput[l](0,5){B} \uput[r](2.6,5){C} \uput[r](1.6,3){D}
	\end{pspicture}
	%Fin PSTricks
\end{minipage}
%\begin{tikzpicture}[scale=1.05]
%% Points basés sur la description : O en bas, A et B alignés ; D et C alignés
%\coordonnée (O) à (0,0);
%\coordonnée (A) à (0,7,3,6);
%\coordonnée (D) à (1,55,3,6);
%\coordonnée (B) à (1.0,5.0);
%\coordonnée (C) à (2.5,5.0);
%
%\dessiner (O) -- (B);
%\dessiner (O) -- (C);
%\dessiner (A) -- (D);
%\dessiner (B) -- (C);
%
%% marques d'angle droit
%\draw ($(A)+(0.18,0)$) -- ($(A)+(0.18,0.18)$) -- ($(A)+(0,0.18)$);
%\draw ($(B)+(0.18,0)$) -- ($(B)+(0.18,0.18)$) -- ($(B)+(0,0.18)$);
%
%\node[gauche] à (O) {$O$};
%\node[gauche] à (A) {$A$};
%\node[right] à (D) {$D$};
%\node[en haut à gauche] à (B) {$B$};
%\node[en haut à droite] à (C) {$C$};
%\end{tikzpicture}

\medskip

\begin{énumérer}
	\item Montrer que la longueur du segment [OA] est égale à 8 cm.
	\item Justifier que les droites (BC) et (AD) sont parallèles.
	\item Calculer la longueur du segment [OB].
	\item Une entreprise souhaite fabriquer des gobelets. Un gobelet (grisé sur le schéma ci-dessous) à la forme d'un tronc de cône (cône coupé par un plan parallèle à sa base).

	\medskip

	\begin{minipage}{0.48\linewidth}
		On reprend les données précédentes :

		\begin{itemize}
			\item O, A, B sont alignés
			\item O, D, C sont alignés
			\item OD = 8,2 cm
			\item AD $= 1,8$ cm
			\item BC $= 4,5$ cm
		\end{itemize}
	\end{minipage}\hfill
	\begin{minipage}{0.48\linewidth}
		\emph{La figure ci-dessous n'est pas à l'échelle.}

		\medskip

		%Début Tikz
% \begin{center}
% \begin{tikzpicture}[scale=0.9]
% % grand cône avec tronc ombré
% \coordonnerie (O) à (0,0);
% \coordinate (Atop) at (-1.0,3.2); % petite ellipse à gauche
% \coordinate (Btop) à (-2.4,4.6); % grande ellipse à gauche
%
% % Lignes du cône externe
% \draw[dashed] (O) -- (-2.4,4.6);
% \draw[dashed] (O) -- (2.4,4.6);
% \draw[dashed] (O) -- (-1.0,3.2);
% \draw[dashed] (O) -- (1.0,3.2);
%
% % Grande ellipse (haut, bord de la tasse)
% \draw (-2.4,4.6) ellipse (2.4 et 0.5);
% % petite ellipse (base de la tasse)
% \draw (-1.0,3.2) ellipse (1.0 et 0.22);
%
% % Côtés du tronc de tronc ombré
% \filldraw[fill=gray!25, draw=black] (-1.0,3.2) -- (-2.4,4.1) arc (180:0:2.4 et 0.5) -- (1.0,3.2) arc(0:180:1.0 et 0.22) -- cycle;
%
% \node[right] à (2.4,4.6) {$C$};
% \node[left] à (-2.4,4.6) {$B$};
% \node[right] à (1.0,3.2) {$D$};
% \node[left] à (-1.0,3.2) {$A$};
% \node[ci-dessous] à (O) {$O$};
%
% \draw[->, épais] (3.6,3.8) -- (1.6,3.9);
% \node[right] at (3.6,3.8) {Le gobelet};
% \end{tikzpicture}
% \end{center}
		%Fin Tikz
		%Début PSTricks
		\begin{pspicture}(-2.6,-0.3)(2.6,5)
			%\psgrid
			†ellipse(0,4.8)(2.3,0.3)
			\psline[linestyle=dashed](-2.34,4.8)(2.34,4.8)
			\psellipticarc(0,2.3)(1.15,0.2){180}{360}
			\psellipticarc[linestyle=dashed](0,2.3)(1.15,0.2){0}{180}
			\psline[linestyle=dashed](0,0)(0,4.8)%OB
			%\psline[linestyle=dashed](-2.3,5)(2.3,5)
			\psline[linestyle=dashed](-1.15,2.3)(1.15,2.3)
			\psline(-2.34,4.8)(-1.13,2.3)
			\psline(2.34,4.8)(1.13,2.3)
			\psframe(0,4.8)(0.25,4.55)\psframe(0,2.3)(0.25,2.55)
			\psline[linestyle=dashed](-1.15,2.3)(0,0)(1.15,2.3)
			\uput[d](0,0){O} \uput[ul](0,2.3){A} \uput[ul](0,4.8){B} \uput[r](2.3,4.8){C} \uput[r](1.15,2.3){D}
			\rput(0,3.6){Le gobelet}
		\end{pspicture}
		%Fin PSTricks
	\end{minipage}


	\begin{center}
		\fbox{\parbox{0.7\textwidth}{\textbf{Rappel :}\\
				Volume d'un cône de révolution :\quad $V = \pi \times R^2 \times H \div 3$,\\ où $R$ est le rayon de la base et $H$ est la hauteur du cône.}}
	\end{center}

	\begin{énumérer}
		\item Calculer le volume du grand cône de hauteur [OB] en cm$^3$, arrondi à l'unité.
		\item Calculer le volume du gobelet, en cm$^3$, arrondi à l'unité.
	\end{énumérer}
\end{énumérer}

\newpage

\section*{Exercice 3 \hfill 4 points}

On considère le \textbf{programme A} suivant :

\begin{center}
	\begin{tikzpicture}[
		box/.style={dessiner, coins arrondis, alignement=centre, largeur minimale=3,6cm, hauteur minimale=0,9cm},
		>={Stealth[longueur=2,5mm]}
		]
		\node[box] (start) at (3,6) {Choisir un nombre};
		\node[box] (carre) at (0,4.5) {Calculer son carré};
		\node[box] (mult7) à (6,4.5) {Multiplicateur par 7};
		\node[box] (mult3) à (0,3) {Multiplicateur par 3};
		\node[box] (add) at (3,1.5) {Additionner les deux nombres};
		\node[box] (sous) à (3,0) {Soustraire 6};

		\draw[->] (début) -- (carré);
		\draw[->] (début) -- (mult7);
		\draw[->] (carré) -- (mult3);
		\draw[->] (mult3) -- (ajouter);
		\draw[->] (mult7) -- (ajouter);
		\draw[->] (ajouter) -- (sous);
	\end{tikzpicture}
\end{center}

\begin{minipage}[t]{0.6\linewidth}
	\begin{énumérer}
	\item Appliquer le \textbf{programme A} au nombre 5.

	\item On utilise un tableau pour trouver les résultats correspondants à quelques nombres comme l'indique le tableau ci-contre.

	\medskip

	Parmi les quatre formules ci-dessous, recopier celle qui a été saisie dans la cellule \textsf{B2}, puis étirée vers le bas afin de calculer les résultats donnés par le programme A.

	\medskip

	\emph{Aucune justification n'est attendue.}

	\begin{tblr}{colspec={XX},hlines, vlines,cells={cmd={\sffamily \small}}}
			= 3 * A2 * 2 + 7 * A2 - 6 & = 3 * 1 * 1 + 7 * 1 - 6 \\
			= 3 * A2 * A2 + 7 * A2 - 6& = 3 * A2 * 2 - 7 * A2 + 6 \\
		\end{tblr}
\end{énumérer}
\end{minipage}
\hfill
\begin{tblr}[t]{hlines={darkgray}, vlines, columns={c, cmd={\sffamily \small}},
			colonne{2,3}={wd=2,7cm,colsep=2pt},
			ligne{1}={rowsep=0pt,ht=7mm, bg=gray9},
			ligne{2-Z}={rowsep=1pt},
			colonne{1}={colsep=0pt,7mm,bg=gray9},
			cellule{1}{1}={gris9},
			cellule{2}{2-3}={cmd={\sffamily \bfseries}}}
			\SetCell{halign=r,valign=b}\tikz[baseline={(0,0.15)}]{\clip(0,0) rectangle (0.6,0.5) ;\fill[gray!50] (0,0)--(0.5,0.5)--(0.5,0);}
			   & A & B \\
			1 & Nombre de départ & Résultat du programme A \\
			2 & -3,5 & 6,25 \\
			3 & -3 & 0 \\
			4 et -2,5 et -4,75 \\
			5 & ​​-2 & -8 \\
			6 & -1,5 & -9,75 \\
			7 & -1 & -10 \\
			8 et -0,5 et -8,75 \\
			9 & 0 & -6 \\
			10 & 0,5 & -1,75 \\
			11 & 1 & 4 \\
			12 & 1,5 & 11,25 \\
			13 & 2 & 20 \\
	\end{tblr}
\begin{enumerate}[start=3]
	\item \`A l'aide du tableau, donner une valeur pour laquelle le \textbf{programme A} donne 0.

	\textit{Aucune justification n'est attendue.}

	\item Si on note $x$ le nombre de départ, donner une expression littérale du programme A en fonction de~$x$.
\end{énumérer}

\newpage

On considère maintenant le \textbf{programme B} suivant :

\begin{center}
	\begin{tikzpicture}[
		box/.style={dessiner, coins arrondis, alignement=centre, largeur minimale=3,6cm, hauteur minimale=0,9cm},
		>={Stealth[longueur=2,5mm]}
		]
		\node[box] (start) at (3,5) {Choisir un nombre};
		\node[box] (mult3) à (0,3.5) {Multiplicateur par 3};
		\node[box] (add3) à (6,2.75) {Additionner 3} ;
		\node[box] (sous2) à (0,2) {Soustraire 2};
		\node[box] (multdeux) at (3,0.5) {Multiplier les deux nombres};

		\draw[->] (début) -- (mult3);
		\draw[->] (début) -- (ajouter3);
		\draw[->] (mult3) -- (sous2);
		\draw[->] (sous2) -- (multdeux);
		\draw[->] (add3) -- (multdeux);
	\end{tikzpicture}
\end{center}

\begin{enumerate}[résumé]
	\item Appliquer le \textbf{programme B} au nombre 5.

	\item Si on note $x$ le nombre de départ, donner une expression littérale du \textbf{programme B} en fonction de~$x$.

	\item Mathis affirme que, quel que soit le nombre qu'il choisit, il trouvera le même résultat avec le \textbf{programme A} et le \textbf{programme B}. A-t-il raison ? Justificateur.

	\item Résoudre l'équation $(3x-2)(x+3) = 0$.

	En déduire les valeurs de $x$ pour lesquelles les programmes A et B donnent 0.
\end{énumérer}

\bigskip

\section*{Exercice 4 \hfill 2,5 points}

\begin{minipage}{0.48\linewidth}
	ABCD est un carré.

	ABF est un triangle équilatéral.

	AF $= 4$ cm.

	Les points F{}, A et E sont alignés.
\end{minipage}\hfill
\begin{minipage}{0.48\linewidth}
	\psset{unité=1cm}
	\begin{pspicture}(-3.5,-1.4)(3.5,4.8)
		%\psgrid
		\psline(-3.2,0)(3.2,0)
		\pspolygon(0;0)(3.2;120)(3.2;180)
		%\rput{30}(0,0){\psframe(0,0)(3.2,3.2)}
		\def\barb{\psline(-0.05,-0.1)(-0.05,0.1)\psline(0.05,-0.1)(0.05,0.1)}
		\rput{30}(0,0){\psline(0;0)(3.2;0)(4.525;45)(3.2;90)}
		\psarc*(0,0){0.6}{0}{30}
		\uput[d](1.6;180){4 cm}
		\uput[d](0,0){A}
		\uput[d](-3.2,0){F} \uput[d](3.2,0){E}\uput[ul](-1.6,2.772){B}\uput[r](3.2;30){D}\uput[ur](1.4,4){C}
		\rput(-1.6,0){\barb}\rput(1.6,0){\barb}
		\rput{30}(1.6;30){\barb}\rput{135}(2,3){\barb}\rput{120}(-0.8,1.4){\barb}
		\rput{600}(-2.4,1.4){\barb}\rput{30}(-0.2,3.6){\barb}
		\psdots[dotstyle=+,dotangle=45,dotscale=1.5](3.2,0)
		\rput(0,-1){\textit{Cette figure n'est pas en vraie grandeur.}}
	\end{pspicture}

\end{minipage}
\begin{énumérer}
	\item Justifier que la mesure de l'angle $\widehat{\text{EAD}}$ est égale à $30^\circ$.

	Dans la suite de l'exercice sur utilisera l'échelle suivante : \textbf{10 pas} dans le programme représente \textbf{1 cm} dans la réalité.

	\item Le \textbf{bloc triangle} ci-dessous permet de tracer un triangle équilatéral et le \textbf{bloc carré} permet de construire un carré :

	\begin{center}
		\begin{tblr}{colspec={|X[7cm,c]|X[7cm,c]|}}
			\hline
			\textbf{Le bloc triangle} & \textbf{Le bloc carré} \\
			\hline
			\begin{scratch}
				\initmoreblocks{définir \namemoreblocks{Triangle}}
				\blockrepeat{répéter \ovalnum{3} fois}
				{
					\blockmove{avancer de \ovalnum{\textbf{J}}}
					\blockmove{tourner \turnleft{} de \ovalnum{\textbf{K}} degrés}
				}
			\end{scratch}
			&
			\begin{scratch}
				\initmoreblocks{définir \namemoreblocks{Carré}}
				\blockrepeat{répéter \ovalnum{4} fois}
				{
					\blockmove{avancer de \ovalnum{\textbf{M}}}
					\blockmove{tourner \turnleft{} de \ovalnum{\textbf{N}} degrés}
				}
			\end{scratch}
			\\ \hline
		\end{tblr}
	\end{center}

	Donner les valeurs de \textbf{J}, \textbf{K}, \textbf{M} et \textbf{N} pour que les blocs triangle et carré permettent de construire un triangle équilatéral de 4 cm de côté et un carré de 4 cm de côté.

	\emph{Aucune justification n'est attendue.}

	\item Le \textbf{programme principal} utilise le bloc \textbf{Triangle} et le bloc \textbf{Carré}.

	L'instruction \og s'orienter vers $90\degres$ \fg{} signifie que l'on s'oriente vers la droite.

	Écrire sur la copie le numéro de la figure obtenue grâce à ce programme.

	\emph{Aucune justification n'est attendue.}
\end{énumérer}

%PSTricks
\begin{tblr}{colspec={XXX}, vline{1,3,4}={solid}, vline{2}={3-4}{solid},
	cellule{4}{1,2}={c},vspan=pair}\hline
	\SetCell[r=1,c=2]{l}\textbf{Figure 1}&&\textbf{Principal du programme}\\
	\SetCell[r=1,c=2]{c} %
	\psset{unité=0,8cm}
	\begin{pspicture}(0,-0.5)(8.5,1.5)
			\def\motif{\pspolygon[linecolor=blue](0;0)(0.75;120)(0.75;180)\psframe[linecolor=blue](0,0)(0.75,0.75)}
			\multido{\n=0,75+1,50}{6}{\rput(\n,0){\motif}}
		\end{pspicture}
	&& \SetCell[r=5,c=1]{c}
		\begin{scratch}[scale=0.85]
		\blockinit{quand \greenflag est cliqué}
		\blockmove{aller à x : \ovalnum{0} y : \ovalnum{-100}}
		\blockmove{s'orienter vers ~\ovalnum{90}}
		\blockpen{effacer tout}
		\blockpen{stylet en position d'écriture}
		\blockrepeat{répéter \ovalnum{6} fois}
		{
			\blockmoreblocks{triangle}
			\blockmove{avancer de \ovalnum{50} pas}
			\blockmove{tourner \turnleft{} de \ovalnum{30} degrés}
			\blockmoreblocks{Carré}
			\blockmove{avancer de \ovalnum{50} pas}
			\blockmove{tourner \turnleft{} de \ovalnum{30} degrés}
		}
	\end{scratch} \\ \cline{1-3}
	\textbf{Figure 2}&\textbf{Figure 3}&\\
		{\psset{unité=0,7cm}
		\begin{pspicture}(-2,-2)(2,2)
			\def\dim{1}%%% modifiable
			\def\motif{\psline[linecolor=blue,linewidth=1.2pt](0,0)(\dim,0)(\dim,-\dim)(0,-\dim)(0,0)(\dim;60)(\dim,0)}
				\multido{\n=30+60,\na=0+60}{6}{\rput{\na}(\dim;\n){\motif}}
		\end{pspicture}}
		&
		{\psset{unité=0,9cm}
		\begin{pspicture}(-2,-2)(2,2)
			\def\motif{\psline[linecolor=blue,linewidth=1.2pt](0,0)(0,-1)(-0.8,-1)(-0.8,-0.2)(0,-0.2)(0,0)}
			\def\motiff{\psline[linecolor=blue,linewidth=1.2pt](0,0)(0.8;120)(0.8;180)(0,0)(1;120)}
			\multido{\n=15+60.0,\na=0+60}{6}{\rput{\na}(1.9;\n){\motif}}
			\multido{\n=15+60.0,\na=0+60}{6}{\rput{\na}(1.9;\n){\motiff}}
		\end{pspicture}} & \\ \cline{1-3}
	\SetCell[c=2]{l}\textbf{Figure 4}&&\\
	\SetCell[c=2]{c}
	{\psset{unité=0,8}
	\begin{pspicture}(0,-0.5)(9,1.5)
			\def\motif{\pspolygon[linecolor=blue](0;0)(0.75;120)(0.75;180)\psframe[linecolor=blue](0,0)(0.75,0.75)\psline(0.75,0)(1,0)}
			\multido{\n=1.00+1.75}{6}{\rput(\n,0){\motif}}
	\end{pspicture}}
	&&\\
	\hline
\end{tblr}
%Fin PSTricks
\end{document}